The \gls{sm} of elementary particles~\cite{Grossman_StandardModel_2023} describes fundamental particles and their interactions, but it does not provide a complete description of nature. For example, it does not explain the nature of dark matter, the origin of neutrino masses, or the matter--antimatter asymmetry of the universe. Many \gls{bsm} theories and models have been developed to address these questions. Examples include:

\begin{itemize}
    \item Supersymmetric models, which introduce a superpartner for each \gls{sm} state and can provide dark matter candidates.
    \item Seesaw models, which introduce right-handed neutrino fields to generate neutrino masses.
    \item Leptogenesis models, which connect neutrino masses and CP violation to the matter--antimatter asymmetry.
\end{itemize}

Particle accelerators are a major tool in searches for \gls{bsm} physics because they combine high collision energies with high luminosities. Examples include the Large Hadron Collider (LHC)~\cite{CERN_Courier_LHC_Run3}, with a collision energy of $13.6\,\mathrm{TeV}$ and an expected total integrated luminosity of $550\,\mathrm{fb}^{-1}$ since its first collisions, and the proposed \gls{fcc-ee}~\cite{Benedikt_2026_FCC}, with energies from $90$ to $365\,\mathrm{GeV}$ and annual integrated luminosities from $0.67$ to $68\,\mathrm{ab}^{-1}$, depending on the interaction point.

Hundreds of searches using LHC data have found no compelling evidence for \gls{bsm} processes; see, for example, Ref.~\cite{PDG_2024} and the references therein. Nevertheless, only a small fraction of possible scenarios has been probed. For example, about $40\%$ of the \gls{pmssm} parameter points excluded by ATLAS evade the available simplified-model constraints~\cite{Ambrogi_2018_pMSSM_misses}. The large, complex datasets collected at colliders therefore remain a valuable resource for further searches.

A conventional search for a specific \gls{bsm} model~\cite{ATLAS_Higgs_2012} involves several labor-intensive steps: identifying production mechanisms and decay modes, defining event selections to enhance the signal relative to the \gls{sm} background, modeling signal and background distributions, quantifying systematic uncertainties, and constructing a statistical model. These steps require substantial time and effort while probing only a limited range of possible scenarios.

In previous work~\cite{Volkovich_2022_DDP, Birman_2022_Symmetries, LastDDPPaper}, we proposed our \gls{ddp} for model-agnostic searches to address these limitations. By scanning detector data for deviations from \gls{sm} predictions, this approach can identify starting points for further investigation. In particular, symmetries of the \gls{sm} impose selection rules and relations between observable processes, providing a reference against which to search for deviations. The symmetrized formalism introduced in Ref.~\cite{LastDDPPaper} provides a statistical method for comparing pairs of data samples.

In this work, we extend that formalism to multiple observables and introduce nuisance parameters to account for systematic effects on the samples.

We apply these extensions to CMS Open Data simulated samples to investigate lepton flavor violation, a well-motivated target for \gls{bsm} searches suited to this approach. We compare pairs of samples selected by lepton flavor that are expected to exhibit the corresponding \gls{sm} symmetry.
